\PassOptionsToPackage{table}{xcolor}
\documentclass[10pt,aspectratio=169]{beamer}
\newcommand{\LectureNumber}{09}
\input{course-theme.tex}
\title{Short-circuiting and switch}
\subtitle{CSC103 Programming Fundamentals / Lecture 09}
\author{Dr. Muhammad Farhan}
\institute{Associate Professor of Computer Science\\Department of Computer Science\\COMSATS University Islamabad\\Sahiwal Campus}
\date{Fall 2026}
\begin{document}
\begin{frame}\titlepage\end{frame}

\begin{frame}[fragile,t]{The problem for this lecture}
\begin{columns}[T,onlytextwidth]\begin{column}{0.60\textwidth}
Write a menu mapping 1 to Add, 2 to Remove and every other integer to Invalid. Test 1, 2 and 9.
\end{column}\begin{column}{0.35\textwidth}\small
Read one integer menu choice.
\end{column}\end{columns}
\note{Introduce the target problem before explaining the tools. Revisit it in the guided solution later.\par Syllabus reading: Deitel Ch 7 (as listed). CLO-2.}
\end{frame}

\begin{frame}[fragile,t]{Learning objectives}
\begin{columns}[T,onlytextwidth]\begin{column}{0.60\textwidth}
\begin{itemize}
\item Use short-circuit guards safely
\item Choose a conditional expression
\item Implement a classic switch statement
\end{itemize}
\end{column}\begin{column}{0.35\textwidth}\small
CLO-2

Short-circuit evaluation\par Conditional operator\par Classic switch selection\par Choice of selection structure
\end{column}\end{columns}
\note{Explain how each objective supports the opening problem. The final questions check these objectives.\par Syllabus reading: Deitel Ch 7 (as listed). CLO-2.}
\end{frame}

\begin{frame}[fragile,t]{Short-circuit evaluation}
\begin{itemize}
\item \&\& skips its right operand when the left operand is false.
\item || skips its right operand when the left operand is true.
\item A guard must precede the operation it protects.
\end{itemize}
\vspace{1mm}\begin{block}{Example}
\begin{lstlisting}
boolean safe = divisor != 0
    && numerator / divisor > 2;
\end{lstlisting}
\end{block}
\note{The division never runs when divisor is zero. Reversing the operands removes that protection. Boolean \& and | evaluate both operands.\par Syllabus reading: Deitel Ch 7 (as listed). CLO-2.}
\end{frame}

\begin{frame}[fragile,t]{Conditional operator}
\begin{itemize}
\item condition ? valueIfTrue : valueIfFalse produces a value.
\item Use it for short alternatives with compatible types.
\item Use if-else for longer branch behaviour.
\end{itemize}
\vspace{1mm}\begin{block}{Example}
\begin{lstlisting}
int larger = a >= b ? a : b;
\end{lstlisting}
\end{block}
\note{Only the selected alternative is evaluated. Numeric type promotion can influence the resulting type, so keep beginner alternatives consistent.\par Syllabus reading: Deitel Ch 7 (as listed). CLO-2.}
\end{frame}

\begin{frame}[fragile,t]{Classic switch selection}
\begin{itemize}
\item switch selects a branch from a selector value.
\item case labels are constant alternatives. default handles unmatched values.
\item Classic colon cases continue into later cases unless control exits.
\end{itemize}
\vspace{1mm}\begin{block}{Example}
\begin{lstlisting}
switch (choice) {
  case 1: label = "Open"; break;
  case 2: label = "Save"; break;
  default: label = "Unknown";
}
\end{lstlisting}
\end{block}
\note{Use the classic statement to align with the outline. Modern arrow cases exist but are optional enrichment. A classic int selector is sufficient here.\par Syllabus reading: Deitel Ch 7 (as listed). CLO-2.}
\end{frame}

\begin{frame}[fragile,t]{Choice of selection structure}
\begin{itemize}
\item if-else handles ranges and compound conditions.
\item switch suits discrete alternatives such as menu choices.
\item Document any intentional fall-through.
\end{itemize}
\vspace{1mm}\begin{block}{Example}
\begin{lstlisting}
Range: temperature < 0
Discrete choice: menu option 1, 2 or 3
\end{lstlisting}
\end{block}
\note{Show that switch does not directly express a relational case such as case score \textgreater{} 50. String switch is possible, but null handling needs care.\par Syllabus reading: Deitel Ch 7 (as listed). CLO-2.}
\end{frame}

\begin{frame}[fragile,t]{A guard protects only later evaluation}
\begin{itemize}
\item Short-circuiting follows left-to-right evaluation.
\item A safe condition first establishes the requirement for the next operation.
\item With integer division, the divisor check must precede the division.
\end{itemize}
\vspace{1mm}\begin{block}{Example}
\begin{lstlisting}
d != 0 && 10 / d > 1
For d=0, the right side never runs.
\end{lstlisting}
\end{block}
\note{Explain the mechanism using the displayed example. Short-circuiting follows left-to-right evaluation. A safe condition first establishes the requirement for the next operation. With integer division, the divisor check must precede the division.\par Syllabus reading: Deitel Ch 7 (as listed). CLO-2.}
\end{frame}

\begin{frame}[fragile,t]{Switch fall-through is an execution rule}
\begin{itemize}
\item A matching classic case determines where execution enters the switch.
\item It does not automatically determine where execution leaves.
\item break transfers control to the statement after the switch.
\end{itemize}
\vspace{1mm}\begin{block}{Example}
\begin{lstlisting}
case 1: print Add; break;
case 2: print Remove; break;
default: print Invalid;
\end{lstlisting}
\end{block}
\note{Explain the mechanism using the displayed example. A matching classic case determines where execution enters the switch. It does not automatically determine where execution leaves. break transfers control to the statement after the switch.\par Syllabus reading: Deitel Ch 7 (as listed). CLO-2.}
\end{frame}


\begin{frame}[fragile,t]{A short-circuit guard controls evaluation}
\begin{center}\begin{tikzpicture}
\node[decision] (a) at (0,0) {d != 0?};
\node[alt] (b) at (-3.6,-1.65) {Evaluate 10 / d \textgreater{} 1};
\node[hot] (c) at (3.6,-1.65) {Result false};
\node[box] (d) at (0,-3.25) {Use Boolean result};
\draw[arr] (a.west) -| node[lab,pos=.25]{true} (b.north);\draw[arr] (a.east) -| node[lab,pos=.25]{false} (c.north);\draw[arr] (b.south) |- (d.west);\draw[arr] (c.south) |- (d.east);
\end{tikzpicture}\end{center}
\begin{block}{What to notice}When d is zero, Java skips the division and the entire AND expression is false.\end{block}
\note{Trace each labelled connection. When d is zero, Java skips the division and the entire AND expression is false.}
\end{frame}
\begin{frame}[fragile,t]{Key distinctions}
\small\renewcommand{\arraystretch}{1.5}
\rowcolors{2}{pale}{white}
\begin{tabularx}{\textwidth}{>{\raggedright\arraybackslash}X>{\raggedright\arraybackslash}X>{\raggedright\arraybackslash}X}
\rowcolor{navy}
\color{white}\bfseries Construct & \color{white}\bfseries Best suited to & \color{white}\bfseries Key behaviour\\
if-else & Ranges and compound rules & First matching branch\\
?: & Short value alternatives & Produces a selected value\\
Classic switch & Discrete menu alternatives & Falls through without an exit\\
\end{tabularx}
\vspace{3mm}\footnotesize These distinctions determine which operation or control structure fits the problem.
\note{\par Syllabus reading: Deitel Ch 7 (as listed). CLO-2.}
\end{frame}

\begin{frame}[fragile,t]{First example: execution}
\begin{lstlisting}
int choice = 2;
String label;
switch (choice) {
  case 1: label = "Open"; break;
  case 2: label = "Save"; break;
  default: label = "Unknown";
}
System.out.println(label);
\end{lstlisting}
\note{This is the main body from Java\_Examples/Lecture09.java. The source file includes the complete class. Predict the result before the next slide.\par Syllabus reading: Deitel Ch 7 (as listed). CLO-2.}
\end{frame}

\begin{frame}[fragile,t]{First example: result}
\begin{columns}[T,onlytextwidth]\begin{column}{0.60\textwidth}
Save\par Execution enters case 2.\par break exits the switch.
\end{column}\begin{column}{0.35\textwidth}\small
Classic switch selection

Choice of selection structure
\end{column}\end{columns}
\note{Expected output and interpretation: Save
Execution enters case 2.
break exits the switch.\par Syllabus reading: Deitel Ch 7 (as listed). CLO-2.}
\end{frame}

\begin{frame}[fragile,t]{Guided practice}
\begin{columns}[T,onlytextwidth]\begin{column}{0.60\textwidth}
Write a menu mapping 1 to Add, 2 to Remove and every other integer to Invalid. Test 1, 2 and 9.
\end{column}\begin{column}{0.35\textwidth}\small
Attempt the solution before continuing.

Use the definitions and examples from this lecture.
\end{column}\end{columns}
\note{Give students 8 minutes to attempt the problem. The following slides provide a complete solution. Original solution guidance: Use two case labels, each ending with break, plus a default branch. Outputs: Add, Remove, Invalid.\par Syllabus reading: Deitel Ch 7 (as listed). CLO-2.}
\end{frame}

\begin{frame}[fragile,t]{Solution strategy}
\begin{columns}[T,onlytextwidth]\begin{column}{0.60\textwidth}
\begin{itemize}
\item Read one integer menu choice.
\item Use distinct case labels for 1 and 2, with a break in each.
\item Use default to report every unmatched integer.
\end{itemize}
\end{column}\begin{column}{0.35\textwidth}\small
Sample console input\par 2\par 
\end{column}\end{columns}
\note{Discuss why each step is needed. State the input assumptions before executing the program.\par Syllabus reading: Deitel Ch 7 (as listed). CLO-2.}
\end{frame}

\begin{frame}[fragile,t]{Complete solution (1/2)}
\begin{lstlisting}
public class Solution09 {
  public static void main(String[] args) throws Exception {
    java.util.Scanner in = new java.util.Scanner(System.in);
    int choice = in.nextInt();
    switch (choice) {
      case 1:
        System.out.println("Add");
        break;
\end{lstlisting}
\note{Complete runnable file: Java\_Examples/Solution09.java. Standard input: 2
  \par Syllabus reading: Deitel Ch 7 (as listed). CLO-2.}
\end{frame}

\begin{frame}[fragile,t]{Complete solution (2/2)}
\begin{lstlisting}
      case 2:
        System.out.println("Remove");
        break;
      default:
        System.out.println("Invalid");
    }
  }
}
\end{lstlisting}
\note{Complete runnable file: Java\_Examples/Solution09.java. Standard input: 2
  \par Syllabus reading: Deitel Ch 7 (as listed). CLO-2.}
\end{frame}

\begin{frame}[fragile,t]{Execution trace}
\small\renewcommand{\arraystretch}{1.5}
\rowcolors{2}{pale}{white}
\begin{tabularx}{\textwidth}{>{\raggedright\arraybackslash}X>{\raggedright\arraybackslash}X>{\raggedright\arraybackslash}X>{\raggedright\arraybackslash}X}
\rowcolor{navy}
\color{white}\bfseries Choice & \color{white}\bfseries Matching case & \color{white}\bfseries Output & \color{white}\bfseries Exit\\
1 & case 1 & Add & break\\
2 & case 2 & Remove & break\\
9 & default & Invalid & End of switch\\
\end{tabularx}
\vspace{3mm}\footnotesize Follow the changing state in execution order.
\note{Manually derived trace for the complete solution. Read one integer menu choice. Use distinct case labels for 1 and 2, with a break in each. Use default to report every unmatched integer.\par Syllabus reading: Deitel Ch 7 (as listed). CLO-2.}
\end{frame}

\begin{frame}[fragile,t]{Solution result}
\begin{columns}[T,onlytextwidth]\begin{column}{0.60\textwidth}
Remove
\end{column}\begin{column}{0.35\textwidth}\small
Why this result follows

Use default to report every unmatched integer.
\end{column}\end{columns}
\note{Expected result: Remove\par Syllabus reading: Deitel Ch 7 (as listed). CLO-2.}
\end{frame}

\begin{frame}[fragile,t]{A common incorrect approach}
boolean ok = 10 / d \textgreater{} 1 \&\& d != 0;
\vspace{1mm}\begin{block}{Example}
\begin{lstlisting}
Put d != 0 first. The current order divides before checking safety and throws when d is zero.
\end{lstlisting}
\end{block}
\note{Ask students to predict the failure before discussing the explanation. The right side gives the correction.\par Syllabus reading: Deitel Ch 7 (as listed). CLO-2.}
\end{frame}

\begin{frame}[fragile,t]{Boundary and failure checks}
\begin{columns}[T,onlytextwidth]\begin{column}{0.60\textwidth}
Check the stated input assumptions.

Create a four-option calculator using switch. Handle an invalid option and guard integer division by zero.
\end{column}\begin{column}{0.35\textwidth}\small
Use two case labels, each ending with break, plus a default branch. Outputs: Add, Remove, Invalid.
\end{column}\end{columns}
\note{Use the specification to derive each expected result. Exercise solution: Use two case labels, each ending with break, plus a default branch. Outputs: Add, Remove, Invalid.\par Syllabus reading: Deitel Ch 7 (as listed). CLO-2.}
\end{frame}

\begin{frame}[fragile,t]{Check your understanding}
\begin{columns}[T,onlytextwidth]\begin{column}{0.60\textwidth}
What does false \&\& riskyCall() do? Why can a missing break change output?
\end{column}\begin{column}{0.35\textwidth}\small
Write a prediction and explain the rule that supports it.
\end{column}\end{columns}
\note{Answer: It skips riskyCall. A missing break in a classic switch can let execution continue into later case bodies.\par Syllabus reading: Deitel Ch 7 (as listed). CLO-2.}
\end{frame}

\begin{frame}[fragile,t]{Answer and explanation}
\begin{columns}[T,onlytextwidth]\begin{column}{0.60\textwidth}
It skips riskyCall. A missing break in a classic switch can let execution continue into later case bodies.
\end{column}\begin{column}{0.35\textwidth}\small
Put d != 0 first. The current order divides before checking safety and throws when d is zero.
\end{column}\end{columns}
\note{Reconnect the answer to the relevant definition rather than treating it as a fact to memorise.\par Syllabus reading: Deitel Ch 7 (as listed). CLO-2.}
\end{frame}


\begin{frame}[fragile,t]{Compare cases and predict outcomes}
\small\renewcommand{\arraystretch}{1.5}\rowcolors{2}{pale}{white}
\begin{tabularx}{\textwidth}{>{\raggedright\arraybackslash}X>{\raggedright\arraybackslash}X>{\raggedright\arraybackslash}X}\rowcolor{navy}
\color{white}\bfseries d & \color{white}\bfseries Right operand evaluated? & \color{white}\bfseries Result\\
0 & No & false\\
2 & Yes: 5 \textgreater{} 1 & true\\
10 & Yes: 1 \textgreater{} 1 & false\\
\end{tabularx}\par\vspace{3mm}\begin{exampleblock}{Discuss}Explain each expected result before running the example. Which case would reveal a wrong assumption?\end{exampleblock}
\note{Read the first two columns and invite predictions before discussing the outcome.}
\end{frame}

\begin{frame}[fragile,t]{Apply the idea}
\begin{alertblock}{Independent practice}Explain why reversing the operands in d != 0 \&\& 10 / d \textgreater{} 1 is unsafe. Write a guarded test and predict it for d=0, 2 and 10.\end{alertblock}\vspace{3mm}\begin{enumerate}\item Work through the input by hand.\item Write or correct the program.\item Justify the expected result.\end{enumerate}\note{Allow 3–5 minutes, then compare answers in pairs. Explain why reversing the operands in d != 0 \&\& 10 / d \textgreater{} 1 is unsafe. Write a guarded test and predict it for d=0, 2 and 10.}
\end{frame}

\begin{frame}[fragile,t]{Solution and reasoning}
\begin{lstlisting}
boolean accepted = d != 0 && 10 / d > 1;
\end{lstlisting}\vspace{1mm}\small\begin{itemize}
\item Java evaluates the left operand first.
\item Putting the division first can throw ArithmeticException before the zero check.
\item The three results are false, true and false.
\end{itemize}\note{Answer: Java evaluates the left operand first. Putting the division first can throw ArithmeticException before the zero check. The three results are false, true and false.}
\end{frame}

\begin{frame}[fragile,t]{Letter grades as a switch problem}
\begin{itemize}
\item A fixed mapping from one letter to one value fits discrete selection.
\item Normalise case once, then reject letters outside the specified grade set.
\item In classic switch syntax, break prevents unintended fall{-}through.
\end{itemize}
\vspace{3mm}\footnotesize\textcolor{accent}{Source: supplied ISB campus slides. See speaker notes and ISB\_Source\_Map.md.}
\note{Adapted from the supplied ISB campus folder: Strings {-}practice.pptx, slides 4, 5; Lecture{-}8{-}9.pptx, slides 16, 17, 18, 20. Re{-}expressed the source if{-}else grade mapping as a classic switch to match this lecture. Source exercise deck credits Instructor Khurram Iqbal.}
\end{frame}

\begin{frame}[fragile,t]{Worked problem: Letter grades as a switch problem}
\begin{block}{Task}Map A, B, C, D and F to 4, 3, 2, 1 and 0. Demonstrate lowercase b after case normalisation.\end{block}
\small Input: Values are fixed in the program for a reproducible demonstration.\par\vspace{3mm}\begin{exampleblock}{Before running}Predict the output and identify the variables that change.\end{exampleblock}
\note{Adapted from the supplied ISB campus folder: Strings {-}practice.pptx, slides 4, 5; Lecture{-}8{-}9.pptx, slides 16, 17, 18, 20. Re{-}expressed the source if{-}else grade mapping as a classic switch to match this lecture. Source exercise deck credits Instructor Khurram Iqbal.}
\end{frame}

\begin{frame}[fragile,t]{Complete ISB example (1/1)}
\begin{lstlisting}
public class IsbLecture09 {
  public static void main(String[] args) throws Exception {
    char grade = Character.toUpperCase('b');
    int points;
    switch (grade) {
      case 'A': points = 4; break;
      case 'B': points = 3; break;
      case 'C': points = 2; break;
      case 'D': points = 1; break;
      case 'F': points = 0; break;
      default: points = -1;
    }
    System.out.println(points);
  }

}
\end{lstlisting}
\note{Adapted from the supplied ISB campus folder: Strings {-}practice.pptx, slides 4, 5; Lecture{-}8{-}9.pptx, slides 16, 17, 18, 20. Re{-}expressed the source if{-}else grade mapping as a classic switch to match this lecture. Source exercise deck credits Instructor Khurram Iqbal. Complete file: ISB\_Java\_Examples/IsbLecture09.java.}
\end{frame}

\begin{frame}[fragile,t]{Worked trace and expected result}
\small\renewcommand{\arraystretch}{1.4}\rowcolors{2}{pale}{white}
\begin{tabularx}{\textwidth}{>{\raggedright\arraybackslash}X>{\raggedright\arraybackslash}X>{\raggedright\arraybackslash}X}\rowcolor{navy}
\color{white}\bfseries Input & \color{white}\bfseries Normalised & \color{white}\bfseries Result\\
b & B & 3\\
F & F & 0\\
E & E & {-}1: invalid grade\\
\end{tabularx}\par\vspace{2mm}
\begin{block}{Expected console output}\ttfamily\footnotesize 3\end{block}
\note{Adapted from the supplied ISB campus folder: Strings {-}practice.pptx, slides 4, 5; Lecture{-}8{-}9.pptx, slides 16, 17, 18, 20. Re{-}expressed the source if{-}else grade mapping as a classic switch to match this lecture. Source exercise deck credits Instructor Khurram Iqbal. Trace and expected values independently worked for this edition.}
\end{frame}

\begin{frame}[fragile,t]{Extend the ISB example}
\begin{alertblock}{Practice}If input is a String rather than a char, what must be checked before charAt(0)?\end{alertblock}\vspace{3mm}\begin{exampleblock}{Explain your reasoning}Require exactly one character; empty input and multi{-}character grades should be rejected.\end{exampleblock}
\note{Adapted from the supplied ISB campus folder: Strings {-}practice.pptx, slides 4, 5; Lecture{-}8{-}9.pptx, slides 16, 17, 18, 20. Re{-}expressed the source if{-}else grade mapping as a classic switch to match this lecture. Source exercise deck credits Instructor Khurram Iqbal. Answer: Require exactly one character; empty input and multi{-}character grades should be rejected.}
\end{frame}
\begin{frame}[fragile,t]{Lecture summary}
\begin{columns}[T,onlytextwidth]\begin{column}{0.60\textwidth}
\begin{itemize}
\item short-circuit guards safely
\item a conditional expression
\item a classic switch statement
\end{itemize}
\end{column}\begin{column}{0.35\textwidth}\small
\begin{itemize}
\item Read one integer menu choice.
\item Use distinct case labels for 1 and 2, with a break in each.
\item Use default to report every unmatched integer.
\end{itemize}
\end{column}\end{columns}
\note{Ask students to give one concrete example for the concepts listed. Use the worked solution to connect the topics.\par Syllabus reading: Deitel Ch 7 (as listed). CLO-2.}
\end{frame}

\begin{frame}[fragile,t]{After-class practice}
\begin{columns}[T,onlytextwidth]\begin{column}{0.60\textwidth}
Create a four-option calculator using switch. Handle an invalid option and guard integer division by zero.
\end{column}\begin{column}{0.35\textwidth}\small
Syllabus reading\par Deitel Ch 7 (as listed)

Next lecture: While loops
\end{column}\end{columns}
\note{Reading labels follow the supplied outline. Check topic headings against the available textbook edition. Require a program and expected results for the practice task.\par Syllabus reading: Deitel Ch 7 (as listed). CLO-2.}
\end{frame}
\end{document}
